\chapter{Print, diagnostyka i czytanie bledow}

\section{Po co istnieje \mc{print(...)}}
Gdy uczysz sie nowego jezyka albo rozwijasz bardziej zlozony projekt, bardzo czesto
chcesz po prostu zobaczyc, co dzieje sie w srodku programu. Wlasnie po to istnieje
funkcja wbudowana \mc{print(...)}.

Dziala ona jak narzedzie podgladu. Pozwala wypisac wartosci do konsoli podczas wykonania
MoonChunk.

\section{Najprostsze uzycie}
\begin{lstlisting}[style=moonchunk,caption={Podstawowe wypisywanie wartosci}]
print("Hello");
print(1, 1.0, true, "text");
\end{lstlisting}

Mozesz przekazac jeden argument albo kilka argumentow jednoczesnie.

\section{Wypisywanie pustej linii}
\mc{print()} bez argumentow tez ma sens. Pozwala odseparowac grupy logow i zrobic
czytelniejszy output w terminalu.

\begin{lstlisting}[style=moonchunk,caption={Pusta linia przez print()}]
print("Sekcja A");
print();
print("Sekcja B");
\end{lstlisting}

\section{Co mozna wypisac}
MoonChunk stara sie byc tutaj elastyczny. Mozesz wypisywac:

\begin{itemize}
  \item liczby,
  \item tekst,
  \item wartosci logiczne,
  \item tablice,
  \item obiekty,
  \item wynik wywolania funkcji,
  \item sama funkcje,
  \item nawet chunki jako obiekty logiczne.
\end{itemize}

\begin{lstlisting}[style=moonchunk,caption={Rozne rodzaje argumentow print}]
function greet(name: string): string {
  return "Hello, " + name;
}

print(greet);
print(greet("MoonChunk"));
print(Main);
\end{lstlisting}

To przydatne szczegolnie przy debugowaniu i poznawaniu modelu wykonania jezyka.

\section{Print jako narzedzie edukacyjne}
Na poczatku nauki \mc{print(...)} jest jednym z najlepszych sposobow zrozumienia programu.
Dzieki niemu mozesz sprawdzac:

\begin{itemize}
  \item czy zmienna ma taka wartosc, jakiej oczekujesz,
  \item czy warunek daje \mc{true} czy \mc{false},
  \item jak zmienia sie licznik w petli,
  \item czy funkcja zwraca poprawny wynik,
  \item jakie dane zostaly wczytane z pliku JSON.
\end{itemize}

\section{Jak wyglada dobry log}
Dobre logowanie nie polega na wypisywaniu wszystkiego bez planu. Lepiej uzyc krotszych,
konkretnych komunikatow.

\begin{lstlisting}[style=moonchunk,caption={Czytelny styl logowania}]
print("counter:", counter);
print("current post slug:", posts[i].slug);
print("isPublished:", posts[i].published);
\end{lstlisting}

Taki zapis pozwala od razu zrozumiec, na co patrzysz.

\section{Bledy skladniowe a bledy wykonania}
MoonChunk moze zglaszac bledy na roznych etapach. Z perspektywy uzytkownika warto
rozrozniac dwie duze kategorie.

\subsection*{Bledy skladniowe}
To problemy takie jak:

\begin{itemize}
  \item brak srednika,
  \item niezamkniety nawias,
  \item nieoczekiwany token,
  \item literowka w konstrukcji jezyka.
\end{itemize}

\subsection*{Bledy wykonania}
To problemy, ktore pojawiaja sie dopiero podczas realnego wykonywania kodu, np.:

\begin{itemize}
  \item nieznana zmienna,
  \item niezgodnosc typu,
  \item dzielenie przez zero,
  \item odczyt z \mc{null},
  \item uzycie \mc{break} poza petla,
  \item brak inicjalizacji \mc{let} przed odczytem.
\end{itemize}

\section{Jak czytac komunikat diagnostyczny}
Dobry komunikat MoonChunk zwykle zawiera:

\begin{itemize}
  \item informacje, co poszlo nie tak,
  \item numer linii,
  \item numer kolumny,
  \item czasem sugestie dotyczaca literowki.
\end{itemize}

Nie czytaj takiego komunikatu jak kary. Traktuj go jak wskazowke. Najczesciej da sie
z niego wyciagnac bardzo konkretna informację o miejscu i naturze problemu.

\section{Typowe komunikaty i ich znaczenie}
\subsection*{Nieznana zmienna}
Jesli MoonChunk mowi, ze zmienna nie jest zadeklarowana, najczesciej oznacza to jedno
z trzech:

\begin{itemize}
  \item literowke w nazwie,
  \item probę uzycia zmiennej poza zasiegiem,
  \item probę użycia zmiennej zanim została utworzona.
\end{itemize}

\subsection*{Zmienna zadeklarowana, ale nie zainicjalizowana}
Ten komunikat pojawia sie, gdy \mc{let} zostal utworzony bez wartosci, a potem odczytany
za wczesnie.

\subsection*{Niezgodnosc typu}
Oznacza, ze interpreter spodziewal sie jednego rodzaju wartosci, a dostal inny.
Najczesciej trzeba poprawic deklaracje typu albo wykonac swiadome rzutowanie.

\subsection*{Blad sterowania przeplywem}
Dotyczy np. \mc{break} albo \mc{continue} poza petla.

\section{Strategia debugowania krok po kroku}
Gdy cos nie dziala, warto przyjac prosty plan:

\begin{enumerate}
  \item znajdz pierwsza linie, na ktorej pojawia sie blad,
  \item sprawdz, jakie wartosci biora w nim udzial,
  \item dodaj \mc{print(...)} przed problematycznym miejscem,
  \item upewnij sie, ze typy sa takie, jak myslisz,
  \item uprosc fragment do mniejszego przykladu i sprawdz go osobno.
\end{enumerate}

To podejscie jest nudne, ale niezwykle skuteczne.

\section{Jak odroznic blad w logice od bledu w danych}
Czasem problem nie lezy w samej skladni programu, tylko w tym, co przychodzi z danych.
Przykladowo:

\begin{itemize}
  \item JSON moze nie miec pola, ktore zakladasz,
  \item liczba moze przyjsc jako tekst,
  \item element listy moze byc pusty,
  \item flaga logiczna moze miec nieoczekiwany stan.
\end{itemize}

W takich sytuacjach \mc{print(...)} jest nieoceniony.

\section{Dobra praktyka: loguj stan, nie emocje}
Zamiast wypisywac rzeczy typu \mc{"ups"} albo \mc{"cos nie dziala"}, lepiej logowac
konkretne stany.

\begin{lstlisting}[style=moonchunk,caption={Przyklad konkretnego logu}]
print("post index:", i);
print("post title:", posts[i].title);
print("published:", posts[i].published);
\end{lstlisting}

Taki log od razu pomaga znalezc przyczyne problemu.

\section{Kiedy usuwac logi}
W projekcie edukacyjnym czesc logow warto zostawic na dluzej. W projekcie produkcyjnym
lepiej usuwac albo ograniczac logi diagnostyczne, gdy przestaja byc potrzebne. Kod jest
wtedy cichszy i latwiejszy do czytania.
